\begin{statement}{E.34}{Proved}{The four-strand order-three torsion class gives a second finite test}
Keep the finite quotient $q:B_4\twoheadrightarrow T$ and the
elements $a,b,c,z,C$ of E.30.  Put
\[
 e=a^2bc,\qquad e_F=a_F^2b_Fc_F
       \quad\text{for }F=(\lambda,f)\in\widehat{\mathrm{GT}}.
 \tag{E.34.1}
\]
Then $e^3=z$, and every full-GT element satisfies
\[
                   e_F C\sim_{T/C} e^\lambda C.
 \tag{E.34.2}
\]
Equivalently there exist $h\in T$ and
$j\in\mathbb Z/|z|$ with
\[
                       e_F=h e^\lambda h^{-1}z^j.
 \tag{E.34.3}
\]
Only the residue of $\lambda$ modulo three enters the
class on the right of \textup{(E.34.2)}.  In particular
both E.24 multiplier residues $1,73$ give the class
$eC$ there.  The test uses the actual homomorphism
$q\circ F$ and the capping map of E.30.
\end{statement}
\begin{proof}
Write
$\delta=\sigma_1\sigma_2\sigma_3$ and
$\epsilon=\sigma_1^2\sigma_2\sigma_3$ in $B_4$.
There is a short proof of the exact braid identity
$\epsilon^3=\delta^4$.  Using a digit $i$ for
$\sigma_i$, its successive positive words are
\[
\begin{aligned}
 112311231123&=112131231123=112131213123\\
             &=121231213123=121232123123\\
             &=121323123123=123123123123.
\end{aligned}
 \tag{E.34.4}
\]
The six changes use respectively distant commutation,
distant commutation, the $121=212$ relation twice,
the $232=323$ relation, and distant commutation.
The program \texttt{e\_periodic\_braid\_root.py}
independently verifies every location and substitution.
Applying $q$ proves $e^3=z$.

Let $\bar\epsilon$ be the image in
$B_4/\langle\delta^4\rangle$, identified with the
fifth-point stabilizer in $\Gamma_{0,[5]}$ by
\textup{(E.30.5)}.  Formula \textup{(E.34.4)} and
its permutation $(234)$ show that its order is
exactly three, with fixed marked points $1,5$.
The actual Nakamura--Schneps action comparison
\textup{(E.30.5a)--(E.30.5c)} makes
$F(\bar\epsilon)$ another element of order three.
The action-independent torsion classification in
Collas's Theorem~3.3 and Corollary~3.2 makes it
conjugate to $\bar\epsilon$ or $\bar\epsilon^2$.
The explicit normalizer below conjugates these
two powers.  Since $\lambda$ is a unit modulo
three, the resulting conjugacy can be written
\[
 F(\bar\epsilon)=u\bar\epsilon^\lambda u^{-1},
                   \qquad u\in\widehat\Gamma_{0,[5]}.
 \tag{E.34.5}
\]
We must locate this conjugator in the actual
fifth-point stabilizer before passing to $T/C$.

There is an explicit geometric normalizer reversing
$\bar\epsilon$ and interchanging its two fixed
points.  The periodic braid in
\textup{(E.34.4)} is the order-three representative
$\epsilon_{n-2}$ for $n=5$ in Collas's
Corollary~3.2 \cite{ERootCollas2012}.
Its rotation model has the five marked points
\[
                    \{0,\infty,1,\zeta_3,\zeta_3^2\},
 \qquad R(t)=\zeta_3t.
 \tag{E.34.6}
\]
The order-two Möbius map $w_0(t)=1/t$ preserves
this set, exchanges $0$ and $\infty$, and satisfies
$w_0Rw_0^{-1}=R^{-1}$ by literal substitution.
Transporting this map by the marked-sphere
identification with the periodic braid gives
$w\in\Gamma_{0,[5]}$ such that
\[
 w\bar\epsilon w^{-1}=\bar\epsilon^{-1},
                 \qquad w(1)=5,\quad w(5)=1.
 \tag{E.34.7}
\]
This is an actual discrete mapping class and hence
also gives an element of its profinite completion.
The two additional equatorial points are exchanged
and the third is fixed, as recorded by the literal
permutation check in the program.

The Drinfeld action preserves the permutation of
$\bar\epsilon$.  Formula \textup{(E.34.5)} therefore
makes $u$ preserve its fixed-point set $\{1,5\}$.
If its permutation exchanges these points, then
$v=uw$ fixes $5$ and \textup{(E.34.7)} gives
\[
                    F(\bar\epsilon)
                       =v\bar\epsilon^{-\lambda}v^{-1}.
 \tag{E.34.8}
\]
Both sides and the conjugator now lie in the
completed fifth-point stabilizer.
The abelianization of that group is $\mathbb Z/12$:
the braid relations identify all three generator
classes, and the killed full twist has exponent
sum twelve.  This also computes its profinite
abelianization, since every finite abelian
quotient has exactly those relations.
In this quotient $[\bar\epsilon]=4$, and the
based GT formulas multiply every braid
abelianization by $\lambda$.
Equation \textup{(E.34.8)} would imply
\[
                     4\lambda=-4\lambda\pmod{12},
                     \quad\text{hence }3\mid\lambda.
 \tag{E.34.9}
\]
The latter contradicts
$\lambda\in\widehat{\mathbb Z}^{\times}$.
Thus $u$ fixes $5$ itself, proving the conjugacy
\textup{(E.34.5)} inside
$\widehat B_4/\overline{\langle\delta^4\rangle}$.

The actual map to $T/C$ factors through this
central quotient.  Its application to
\textup{(E.34.5)} proves \textup{(E.34.2)}, and
lifting the conjugator to $T$ proves
\textup{(E.34.3)}.  Since $e^3=z$, its class in
$T/C$ has order dividing three, which proves
the claimed multiplier scope.  This argument
requires no invariance of the finite kernel
of $q$ under the GT action.  The earlier
order-two test and this order-three test use
their respective explicit periodic elements
and the same equivariant capping homomorphism.
\end{proof}
